\documentclass{article}
\usepackage[UTF8]{ctex}
\usepackage{amsmath,amssymb,amsfonts}
\usepackage{enumitem}
\usepackage{geometry}
\geometry{a4paper, margin=1in}
% === 错题本排版优化（不参与题目标识） ===
\setlength{\parindent}{0pt}
\setlength{\parskip}{0.35em plus 0.10em minus 0.08em}
\setlength{\emergencystretch}{3em}
\setlist{itemsep=0.25em, topsep=0.40em, parsep=0pt}
\allowdisplaybreaks[2]
% === 错题本排版优化结束 ===
\begin{document}
\section*{浙江大学2025年数学分析真题}
\section*{102．浙江大学2025年数学分析试卷}

\subsection{题目1}
1． \(\lim _{n \rightarrow \infty} n^{2}\left(\arctan \frac{2024}{n}-\arctan \frac{2024}{n+1}\right)\) ．

\subsection{题目2}
2．\(z=z(x, y)\) 由 \(f(z-x, z-y)\) 确定，\(f(u, v)\) 有一阶偏导数，且 \(\frac{\partial f}{\partial u}+\frac{\partial f}{\partial v} \neq 0\) ，计算 \(\frac{\partial z}{\partial x}+\frac{\partial z}{\partial y}\) ．

\subsection{题目3}
3．计算 \(\int_{-8}^{8} \frac{1+\sin x}{1+\sqrt[3]{x^{2}}} \mathrm{~d} x\) ．

\subsection{题目4}
4．\(\sum=\left\{(x, y, z) \mid x^{2}+y^{2}+z^{2}=1, z \geq 0\right\}\) ，取上侧，计算

\[
\iint_{\Sigma} x^{2} \mathrm{~d} y \mathrm{~d} z+y^{2} \mathrm{~d} z \mathrm{~d} x+\left[z^{2}+\sin (x y)\right] \mathrm{d} x \mathrm{~d} y .
\]

二．设 \(\left\{a_{n}\right\}\) 为非负递减数列，且级数 \(\sum_{n=1}^{\infty} a_{n}\) 收敛，证明： \(\lim _{n \rightarrow \infty} n a_{n}=0\) 。

三．设 \(f(x)\) 在闭区间 \([0,1]\) 上连续，且对任意的 \(x\in(0,1)\)，\(f_{+}^{\prime}(x)\) 存在．若 \(f(0)=f(1)\) ，证明：存在\(\xi \in(0,1)\) ，设得 \(f_{+}^{\prime}(\xi) \leq 0\) ．

四．叙述闭区间套定理，用该定理证明闭区间上连续函数可达最大值．

五．设

\[
f(x, y)= \begin{cases}(x+y)^{2} \sin \frac{1}{x^{2}+y^{2}}, & (x, y) \neq(0,0) ; \\ 0, & (x, y)=(0,0) .\end{cases}
\]

讨论 \(f\) 和 \(f\) 的偏导数在 \((0,0)\) 处的连续性及 \(f\) 在 \((0,0)\) 的可微性．

六．证明：\(\sum_{n=1}^{\infty} n^{2} e^{-n x}\) 在 \((0,+\infty)\) 上非一致收敛，在 \([1,+\infty)\) 上一致收敛。

七．证明：\(f(x)\) 在 \((0,1)\) 上一致连续的充要条件是对任意 \((0,1)\) 上的 Cauchy 列 \(\left(x_{n}\right),\left\{f\left(x_{n}\right)\right\}\) 也为Cauchy 列。

八．设 \(\left\{r_{n}\right\}_{n=1}^{\infty}\) 为全体有理数，令

\[
f(x)=\sum_{n: r_{n}<x} \frac{1}{2^{n}} .
\]

其中上述和式是对所有满足 \(r_{\mathrm{n}}<x\) 的正整数 \(n\) 求和。证明：\(f(\hat{x})\) 在无理点连续，在有理点处间断，且单调递增。

九．对任意的实数 \(x\) 与正整数 \(n\) ，记 \(\varphi_{n}(x)=\frac{n}{\sqrt{\pi}} e^{-n^{2} x^{2}}\) ．证明：对任意在实轴上连续且有界的函数 \(f\) ，对任意的实数 \(x\) ，有

\[
\lim _{n \rightarrow \infty} \int_{-\infty}^{+\infty} f(y) \varphi_{n}(x-y) \mathrm{d} y=f(x)
\]
\end{document}
